\documentclass[10pt,a4paper]{amsart}
\usepackage[margin=19mm]{geometry}
\usepackage{amsmath,amssymb,mathtools}
\usepackage[scr=rsfs,cal=euler]{mathalfa}
\usepackage{xfrac}
\usepackage[unicode,hidelinks,bookmarks=false]{hyperref}
\newcommand{\ie}{i.e.\ }
\newcommand{\cf}{cf.\ }
\newcommand{\resp}{respectively\ }
\newcommand{\Subsection}{\subsection}
\providecommand{\nobreakdash}{\nobreak\hbox{-}}
\setlength{\emergencystretch}{2em}
\multlinegap0pt
\usepackage{mathbbol}
\usepackage{setspace}
\usepackage{parskip}
\allowdisplaybreaks
\hypersetup{colorlinks=true,
linkcolor=blue,
citecolor=blue,
urlcolor=cyan}
\newtheorem{theorem}{Theorem}[section]
\newtheorem{lemma}{Lemma}[section]
\newtheorem{remark}{Remark}[section]
\newtheorem{conjecture}{Conjecture}[section]
\begin{document}
\title[An Integral Mean Value Theorem for Weyl Sums over Broken Arc]{An Integral Mean Value Theorem for Weyl Sums over Broken Arcs}
\author{YaoJie Guo}
\date{}
\maketitle
\noindent\emph{Source and licence.} An Integral Mean Value Theorem for Weyl Sums over Broken Arcs. Version arXiv:2608.04948v1 (CC BY 4.0). This material is distributed under the Creative Commons Attribution 4.0 International licence, http://creativecommons.org/licenses/by/4.0/, as recorded in the arXiv OAI licence metadata for this exact version and shown on the abstract page https://arxiv.org/abs/2608.04948v1. Full rights notice: licenses/arxiv-2608.04948-CC-BY-4.0.txt.\par\smallskip
\noindent\textit{Editorial scope.} Counted unit: the original \texttt{lemma} environment \texttt{-} at source lines 133--143 together with its complete proof at lines 144--204; the delivered document keeps source lines 131--220 so that every referenced label and every citation resolves inside it. Native editable source is the paper's own concatenated TeX; the AMS wrapper is editorial formatting only. No publisher class, upstream script or unrelated artwork is redistributed.
\par\smallskip\noindent\textit{Reference note.} This article ships no compiled bibliography (biblatex); the entries delivered here are composed from the author's own BibTeX (.bib) fields, with no field invented and no entry dropped.
\section{The pointwise bound according to Vinogradov's main value theorem}
In this section we will prove the following lemmas:(By the way the definition of letter $m$ in this section in not the same as it in the next section.)

\begin{lemma}
We set $f(x)=\sum_{1\leq i\leq k+1}\alpha c_{i}x^i$, where $l,N,X,Y$ are four integers, where $Y<X$ is sufficiently large and $X=o(N)$, $l>k(k+1)/2$. In this condition, for any given $\epsilon>0$, we have
\begin{equation}
|\sum\limits_{n=N}^{N+X}e(f(n))|\ll Y+X^{1-\frac{kc}{4l}+\epsilon}
\end{equation}
hold true for almost any given $\alpha$ belong to
\[
\mathfrak{m}(Y,X)=\bigcap_{j\leq k-1}\left\{\alpha:\forall q<Y^j,h<q,(h,q)=1,\Big|\alpha-\frac{h}{q}\Big|>\frac{1}{X^{k-j-c}q}\right\}
\]
and $\ll$ is a absolute constant.
\end{lemma}

\begin{proof}
We set $g(n)=f(n+N)$, and we know
\[
|\sum\limits_{n=N}^{N+X}e(f(n))|=|\sum\limits_{n=0}^{X}e(g(n))|
\]
and the first coefficient of $g(n)$ is the same as $f(n)$, so we just need to prove
\[
|\sum\limits_{n=0}^{X}e(f(n))|\ll Y+X^{1-\frac{kc}{4l}+\epsilon}
\]
From \cite{PP1986}, we know 
\[
|\sum\limits_{n=0}^{X}e(f(n))|=Y^{-1}\sum\limits_{y\leq Y}\sum\limits_{n\leq X}e(f(x+y))+O(Y)
\]
and we have
\[
Y^{-1}\Big|\sum\limits_{y\leq Y}\sum\limits_{n\leq X}e(f(x+y))\Big|^{2l}\leq Y^{-1}X^lI^{1/2}(X;k,l)\Big(\sum\limits_{\lambda_1,...\lambda_k}\big|\sum\limits_{y\leq Y}e(f^{(1)}(y)\lambda_1+...+\frac{1}{k!}f^{(k)}(y)\lambda_k)\big|^2\Big)^{1/2}
\]
where $\lambda_j<lX^j$ and
\[
I(X;k,l)=\int_{[0,1]^k}\bigg|\sum\limits_{u\leq X}e(a_ku^k+...+a_1u)\bigg|^{2l}\mathrm{d}\textbf{a}
\]
By Vinogradov's main conjecture(To see in \cite{bourgain2016proofmainconjecturevinogradovs} for $k>3$, and \cite{wooley2014cubiccasemainconjecture} in condition that $k=3$) we know
\[
I(X;k,l)\ll X^{\epsilon}(X^{l}+X^{2l-k(k+1)/2})
\]
From Holder inequality(The detail can be found in \cite{PP1986}, just some elementary skills), we have
\[
Y^{-2l}\Big|\sum\limits_{y\leq Y}\sum\limits_{n\leq X}e(f(x+y))\Big|^{2l}\leq Y^{-1}X^{2l-k(k+1)/4+\epsilon}\Big(\sum\limits_{\lambda_1,...\lambda_k}\big|\sum\limits_{y\leq Y}e(f^{(1)}(y)\lambda_1+...+\frac{1}{k!}f^{(k)}(y)\lambda_k)\big|^2\Big)^{1/2}
\]
Then we know
\begin{align*}
&\sum\limits_{\lambda_1,...\lambda_k}\big|\sum\limits_{y\leq Y}e(f^{(1)}(y)\lambda_1+...+\frac{1}{k}f^{(k)}(y)\lambda_k)\big|^2\\
=&\sum\limits_{y\leq Y}\sum\limits_{y'\leq Y}\sum\limits_{\lambda_1,...\lambda_k}e((f^{(1)}(y)-f^{(1)}(y'))\lambda_1+...+\frac{1}{k!}(f^{(k)}(y)-f^{(k)}(y'))\lambda_k)\\
=&\sum\limits_{y\leq Y}\sum\limits_{y'\leq Y}\prod\limits_{1\leq j\leq k}\sum\limits_{\lambda_j<lX^j}e\big(\frac{1}{j!}(f^{(j)}(y)-f^{(j)}(y'))\big)
\end{align*}
We know
\[
\sum\limits_{\lambda_j<kX^j}e\big(\frac{1}{j!}(f^{(j)}(y)-f^{(j)}(y'))\lambda_j\big)\ll_k\min\bigg(X^j,\frac{1}{2\|\frac{1}{j!}(f^{(j)}(y)-f^{(j)}(y'))\|}\bigg)
\]
From the definition of $\mathfrak{m}(Y,P)$, we know
\[
\|\frac{1}{j!}(f^{(j)}(y)-f^{(j)}(y'))\|\gg X^{j}
\]
Then we know
\[
\prod\limits_{1\leq j\leq k}\sum\limits_{\lambda_j<lX^j}e\big(\frac{1}{j!}(f^{(j)}(y)-f^{(j)}(y'))\big)\ll\prod\limits_{1\leq j\leq k}\frac{1}{\|\frac{1}{j!}(f^{(j)}(y)-f^{(j)}(y'))\|}
\]
Then we have
\[
\prod\limits_{1\leq j\leq k}\sum\limits_{\lambda_j<lX^j}e\big(\frac{1}{j!}(f^{(j)}(y)-f^{(j)}(y'))\big)\ll X^{(k+1)k/2-kc}
\]
and
\[
\sum\limits_{\lambda_1,...\lambda_k}\big|\sum\limits_{y\leq Y}e(f^{(1)}(y)\lambda_1+...+\frac{1}{k!}f^{(k)}(y)\lambda_k)\big|^2\ll Y^2X^{k(k+1)/2-kc}
\]
Hence we have
\[
Y^{-1}\sum\limits_{y\leq Y}\sum\limits_{n\leq X}e(f(x+y))\ll X^{2l-\frac{kc}{2}+\epsilon}
\]
So we complete the proof.
\end{proof}

To ensure our estimate is effective for nearly almost any given $\alpha$, we need
\[
\mathfrak{L}([0,1]-\mathfrak{m}(Y,X))=o(1)
\]
which means
\[
\frac{Y^j}{X^{k-j-c}}=o(1),\forall j>1
\]
choosing $Y=(\log x)^A$, $c=1/2+\epsilon$ and we have following theorem:
\begin{lemma}
Let $f(x)\in\textbf{Z}[x],\deg f=k$ is a monic polynomial, and $H=[N-N^\theta,N+N^\theta]$. The following estimate
\[
\sum\limits_{\substack{x\in H}}e(\alpha f(x))\ll_\epsilon N^{\theta(1-\frac{1}{4(k+1)})+\epsilon}
\]
hold true for any given $\epsilon>0$ and any $\alpha\in\mathfrak{m}$.
\end{lemma}

\begin{thebibliography}{99}
\bibitem[]{PP1986} Pan,C.D.; Pan,C.B.. The fundamental of Analytic Number Theory. (1986)
\bibitem[]{bourgain2016proofmainconjecturevinogradovs} Jean Bourgain; Ciprian Demeter; Larry Guth. Proof of the main conjecture in Vinogradov's mean value theorem for degrees higher than three. Ann.Math. (2016) 633-682(2)
\bibitem[]{wooley2014cubiccasemainconjecture} Trevor D. Wooley. The cubic case of the main conjecture in Vinogradov's mean value theorem. Adv.Math. (2016) 532-561
\end{thebibliography}
\end{document}
